\documentclass[11pt,a4paper,onecolumn]{article}
\usepackage[spanish,es-noshorthands]{babel}
\usepackage[utf8]{inputenc}
\usepackage[T1]{fontenc}
\usepackage{amsmath,amssymb,amsbsy,amsopn}
\usepackage{graphicx}
\usepackage{booktabs}
\usepackage{array}
\usepackage{multirow}
\usepackage{xcolor}
\usepackage{hyperref}
\usepackage{url}
\usepackage[margin=2.5cm]{geometry}
\usepackage{setspace}
\usepackage{natbib}
\usepackage{fancyhdr}
\onehalfspacing
\pagestyle{fancy}
\fancyhf{}
\fancyhead[L]{Dark Forest Monte Carlo}
\fancyhead[R]{Juan Galaz}
\fancyfoot[C]{\thepage}
\renewcommand{\headrulewidth}{0.4pt}
\newcommand{\doi}[1]{\href{https://doi.org/#1}{#1}}

\title{\bfseries Del Silencio a la Concentración: Transición de Régimen en la Hipótesis del Bosque Oscuro Revelada por Simulaciones de Control con Dinámica de Destrucción Desactivada}
\author{Juan Galaz\\[2mm]
\small ORCID: \href{https://orcid.org/0009-0007-7474-7560}{0009-0007-7474-7560}\\
\small DOI: \href{https://doi.org/10.5281/zenodo.20451755}{10.5281/zenodo.20451755}\\
\small Santiago, Chile\\
\small \texttt{juan.galaz@proton.me}}
\date{10 de mayo de 2026}

\begin{document}
\maketitle

\begin{abstract}
Presentamos una simulación Monte Carlo estocástica de la hipótesis del
Bosque Oscuro, la conjetura de que las civilizaciones galácticas se
destruyen mutuamente de forma preventiva al detectarse, y
caracterizamos las tasas de destrucción resultantes y las distribuciones
espaciales en dos regímenes de población.
Bajo parámetros de Drake \emph{pesimistas} ($N_c = 440$, 50 semillas
independientes) el modelo produce una tasa de destrucción de
$16{,}1 \pm 1{,}6$ por ciento, con civilizaciones supervivientes
distribuidas a una distancia media al vecino más cercano de
$1837 \pm 52$\,al; una prueba de Kolmogórov--Smirnov
confirma que esta distribución es significativamente superpoissoniana
($p = 4{,}89\times10^{-10}$, semilla de referencia,
$\text{cociente}_\text{obs/Poisson} = 1{,}43 \pm 0{,}05$),
consistente con la concentración geométrica de nacimientos en la ZGH.
Simulaciones de control con $P_\text{ataque} = 0$ (10 semillas,
subconjunto apareado de las 50 principales) producen un cociente de control
$r_\text{obs/P}^\text{ctrl} = 1{,}42 \pm 0{,}04$, indistinguible
del cociente con dinámica DF activa
($\Delta r = +0{,}007 \pm 0{,}066$, $0{,}1\sigma$), indicando que la
estructura espacial de los supervivientes en el régimen de baja densidad
está dominada por el perfil ZGH, no por la dinámica de destrucción.
Bajo parámetros de Drake \emph{optimistas} ($N_c = 10{,}000$ acotado,
10 semillas) la tasa de destrucción asciende a $85{,}8 \pm 0{,}2$ por
ciento, con supervivientes concentrados a $478 \pm 17$\,al y el cociente
obs/Poisson se \emph{invierte} a $0{,}578 \pm 0{,}020$ (subpoissoniano).
Simulaciones de control en este régimen ($n = 3$ semillas, viable gracias
a la vectorización del algoritmo de detección) producen
$r_\text{obs/P}^\text{ctrl} = 1{,}117 \pm 0{,}001$, revelando un efecto
diferencial masivo del Bosque Oscuro:
$\Delta r_\text{opt} = -0{,}538 \pm 0{,}020$ ($27\sigma$).
La transición es completa: en baja densidad la dinámica DF es
indistinguible de la geometría ZGH ($\Delta r \approx 0$), mientras que
en alta densidad el DF concentra los supervivientes de forma abrumadora,
multiplicando el agrupamiento espacial por encima del nivel ZGH.
Un análisis de sensibilidad de Sobol global ($N = 4{,}096$,
$49{,}152$ evaluaciones, 5 parámetros) produce índices de primer orden
cercanos a cero ($S_1 \in [-0{,}02,\;0{,}07]$) con efectos totales
$S_T < 1{,}0$ en todos los parámetros (convergencia alcanzada,
intervalos de confianza $\pm 0{,}04$), confirmando que ningún parámetro
domina individualmente y que las interacciones de alto orden controlan la
varianza de la tasa de destrucción.
Estos hallazgos establecen una transición completa dependiente de la
densidad: en el régimen de baja densidad, la dinámica de destrucción es
indistinguible de la geometría ZGH ($\Delta r \approx 0$); en el régimen
de alta densidad, la dinámica DF domina de forma abrumadora la estructura
espacial ($\Delta r = -0{,}54$, $27\sigma$), concentrando a los
supervivientes muy por encima del nivel de agrupamiento geométrico.
\end{abstract}

\vspace{2mm}
\noindent\textbf{Palabras clave:} astrobiología --- inteligencia extraterrestre --- paradoja de Fermi ---
métodos Monte Carlo --- estadística espacial --- zona galáctica habitable
\vspace{5mm}

\section{Introducción}
\label{sec:intro}
% ─────────────────────────────────────────────────────────────────────────────

La paradoja de Fermi, la aparente contradicción entre la alta
probabilidad a priori de civilizaciones tecnológicas extraterrestres y
su ausencia observacional, ha motivado una amplia clase de hipótesis
explicativas \citep{fermi1950,hart1975,brin1983,webb2002}.
Entre las más radicales se encuentra la \emph{hipótesis del Bosque
Oscuro} (HBO), en la que el silencio del cosmos es un equilibrio
estratégico: cualquier civilización que revela su ubicación arriesga
la aniquilación por un vecino más avanzado, por lo que la estrategia
óptima es el silencio o el ataque preventivo
\citep{cixin2008,devladar2013}.

Los modelos formales de la HBO siguen siendo escasos en la literatura
astrobiológica revisada por pares.
\citet{forgan2009} desarrolló el banco de pruebas numérico más citado
para hipótesis de la paradoja de Fermi, utilizando realizaciones Monte
Carlo de la ecuación de Drake, pero modeló la destrucción como una
probabilidad fija ($P_\text{destroy} = 0{,}5$) independiente de la
detección, la distancia o la disparidad tecnológica.
\citet{devladar2013} formalizó el fundamento de teoría de juegos de la
HBO, demostrando que el ataque preventivo es un equilibrio de Nash bajo
ciertas matrices de pagos, pero no simuló las consecuencias espaciales
a escala galáctica.
\citet{foucher2023} y \citet{forgan2011} señalaron independientemente
que la distribución espacial de los supervivientes debería depender de
la densidad de población, pero ninguno cuantificó la transición de
régimen.

El presente trabajo aborda esta brecha implementando un pipeline
Monte Carlo completo de detección--decisión--destrucción con los
siguientes elementos novedosos:
\begin{enumerate}
  \item Una función de probabilidad de detección tridimensional
        continua basada en la edad tecnológica y un radio máximo de
        detección físicamente motivado
        \citep[$r_\text{max} = 1000$\,al,][]{troitskij1989,loeb2007};
  \item Tipos sociológicos heterogéneos (agresivo, defensivo, pacifista,
        curioso) con umbrales de ataque dependientes del tipo, siguiendo
        a \citet{devladar2013};
  \item Un perfil de densidad continuo de la Zona Galáctica Habitable
        (ZGH) basado en \citet{lineweaver2004}, en reemplazo de la
        aproximación de anillo uniforme;
  \item Una memoria histórica de civilizaciones destruidas que reduce la
        probabilidad de comunicación en regiones donde se han observado
        destrucciones previas;
  \item Análisis de doble escenario (parámetros de Drake
        pesimistas/optimistas) que permite la comparación de regímenes;
  \item Análisis de sensibilidad de Sobol global con $N = 4{,}096$
        muestras de Saltelli \citep{saltelli2010}.
\end{enumerate}

El resultado central es una inversión dependiente de la densidad del
cociente espacial obs/Poisson, consistentemente observada en ambos
escenarios.
Para aislar si esta inversión es un efecto genuino de la dinámica DF
o una consecuencia de la concentración geométrica de la ZGH, el trabajo
incluye simulaciones de control con dinámica de destrucción desactivada
(Sección~\ref{sec:control}), cuyo contraste apareado permite cuantificar
el efecto diferencial neto del Bosque Oscuro ($\Delta r$).

Este modelo debe entenderse como una exploración del espacio de
parámetros de la HBO, no como una predicción cuantitativa del estado
actual de la galaxia.
Cinco parámetros carecen de restricción observacional directa
($\theta_\text{threat}$, $d_0$, las fracciones sociológicas,
$p_\text{fn}$, $p_\text{fp}$); los valores absolutos de las tasas de
destrucción dependen de ellos, pero el análisis de sensibilidad global
(\S\ref{sec:sobol}) demuestra que la inversión cualitativa del cociente
obs/Poisson entre regímenes de densidad es robusta ante variaciones
de un factor $3$--$5\times$ en cada parámetro.
El código fuente y los datos de salida son depositados en un repositorio
público para facilitar la reproducción y extensión del modelo.

% ─────────────────────────────────────────────────────────────────────────────
\section{Modelo}
\label{sec:model}
% ─────────────────────────────────────────────────────────────────────────────

\subsection{Generación de civilizaciones}
\label{sec:civgen}

Las civilizaciones se generan muestreando la ecuación de Drake:
\begin{equation}
  N_c = N_\star \cdot f_p \cdot n_e \cdot f_l \cdot f_i \cdot f_c
  \label{eq:drake}
\end{equation}
donde $N_\star = 2\times10^{11}$ es el número de estrellas de la Vía
Láctea aptas para planetas con vida, $f_p = 1{,}0$ (fracción con
planetas, caso conservador que maximiza la población dado el producto
de filtros biológicos), $n_e = 0{,}22$ el número de planetas en zona
habitable por sistema \citep{fressin2013}, produciendo
$N_\text{hp} = N_\star \cdot f_p \cdot n_e = 4{,}4\times10^{10}$
planetas habitables potenciales, y $f_l$, $f_i$, $f_c$ son las
fracciones de vida, inteligencia y comunicación cuyo producto define el
escenario (Tabla \ref{tab:scenarios}).

Las posiciones galácticas se extraen del perfil de densidad continuo de
la ZGH \citep{lineweaver2004}:
\begin{equation}
  f(r) \propto \exp\!\left[-\frac{(r - r_\text{peak})^2}{2\sigma_r^2}\right],
  \quad r_\text{peak} = 25{,}000\,\text{al},\quad \sigma_r = 8{,}000\,\text{al}
  \label{eq:ghz}
\end{equation}
con una semialtura del disco de 300\,al muestreada de una distribución
exponencial con signo aleatorio.
Los tiempos de nacimiento se extraen uniformemente sobre la edad
estimada de la galaxia de $1{,}36\times10^{10}$\,años, y los tiempos
de vida tecnológicos siguen una distribución log-normal con mediana
$10^4$\,años \citep{sandberg2018}.

\begin{figure}
  \centering
  \includegraphics[width=\columnwidth]{output_analisis/analisis_ghz.png}
  \caption{Distribución radial de las civilizaciones generadas respecto al
           centro galáctico (seed\,=\,42, $N_c = 440$).
           Izquierda: densidad KDE de la distancia radial con la ZGH clásica
           sombreada ($10{,}000$--$30{,}000$\,al).
           Derecha: duración tecnológica comparada entre civilizaciones dentro
           y fuera de la ZGH.
           El $42{,}95\%$ de las civilizaciones se ubica dentro del anillo
           clásico de la ZGH, reflejando el perfil continuo de Lineweaver 2004
           (Ecuación~\ref{eq:ghz}) que distribuye nacimientos más allá de los
           límites convencionales.}
  \label{fig:ghz}
\end{figure}

\subsection{Heterogeneidad sociológica}
\label{sec:socio}

A cada civilización se le asigna uno de cuatro tipos sociológicos:
\emph{agresivo} (probabilidad 0{,}30), \emph{defensivo} (0{,}40),
\emph{pacifista} (0{,}20), \emph{curioso} (0{,}10).
En cada paso temporal, cada civilización adopta una de tres
estrategias activas -- \textsc{atacar}, \textsc{esconderse} o
\textsc{comunicar} -- muestreada de un vector de probabilidades
dependiente del tipo.
Las probabilidades base [\textsc{atacar}/\textsc{esconderse}/\textsc{comunicar}]
son: agresivo (0{,}65/0{,}10/0{,}25), defensivo (0{,}20/0{,}60/0{,}20),
pacifista (0{,}05/0{,}75/0{,}20) y curioso (0{,}10/0{,}15/0{,}75).
La estrategia modifica la detectabilidad de la civilización: como
\emph{presa}, los modificadores sobre $P_\text{det}$ son
$\times 1{,}20$ (\textsc{atacar}), $\times 0{,}001$
(\textsc{esconderse}) y $\times 3{,}00$ (\textsc{comunicar}); como
\emph{cazadora}, el modificador es $\times 1{,}20$ solo para
\textsc{atacar}.
La decisión de ataque final sigue una regla probabilística: la
probabilidad base de atacar es 0{,}90 (agresivo), 0{,}60 (defensivo),
0{,}10 (pacifista) y 0{,}30 (curioso), multiplicada por 1{,}50 si la
disparidad tecnológica relativa de la presa supera $\theta_\text{threat}$.

Las probabilidades base de estrategia se ajustan dinámicamente en
función de dos factores contextuales, ambos calculados a partir de
posiciones destruidas observadas localmente -- sin necesidad de
comunicación entre civilizaciones, preservando el silencio del Bosque
Oscuro:
\begin{enumerate}
  \item \emph{Presión ambiental inmediata} (radio de 500\,al):
        si hay $k$ civilizaciones destruidas dentro de 500\,al,
        la probabilidad de \textsc{esconderse} aumenta en
        $\delta = \min(0{,}40,\; 0{,}10\,k)$, mientras que
        \textsc{atacar} y \textsc{comunicar} se reducen cada una en
        $\delta/2$, con renormalización del vector resultante.
  \item \emph{Memoria histórica de zonas peligrosas} (radio de 2{,}000\,al):
        si existe al menos una civilización destruida dentro de un radio
        esférico de 2{,}000\,al, la probabilidad de \textsc{comunicar}
        se reduce en un 80 por ciento, y la masa redistribuida se
        transfiere íntegramente a \textsc{esconderse}, seguido de
        renormalización.
        El radio de memoria (2{,}000\,al) es mayor que el de presión
        inmediata (500\,al) para capturar el aprendizaje de amenazas
        en la vecindad galáctica amplia.
\end{enumerate}

\subsection{Probabilidad de detección}
\label{sec:detection}

La probabilidad de que la civilización $i$ detecte a la civilización
$j$ a distancia euclidiana $d$ en el tiempo $t$ es:
\begin{equation}
  P_\text{det}(d, \tau_i, \nu_i) =
    p_\text{fp} + (1 - p_\text{fn} - p_\text{fp})
    \cdot \frac{\rho(\tau_i,\nu_i)}{1 + (d/d_0)^2}
  \label{eq:pdet}
\end{equation}
donde $\rho(\tau_i,\nu_i) = \min[1,\, r_\text{eff}(\tau_i,\nu_i)/r_\text{max}]$
es la capacidad tecnológica relativa del cazador, $d_0 = 100$\,al es
la escala de atenuación de la señal —correspondiente al alcance de
detección de fugas de radar con telescopios de próxima generación (SKA).
Para una galaxia como M31, con aproximadamente $10^4$ fuentes de radio
detectables en el anillo de habitabilidad, un cociente $r_\text{obs}/P
\approx 0{,}6$ sería distinguible de $1{,}0$ con una prueba KS de
potencia $> 0{,}9$ ($n = 10^4$, tamaño de efecto $d = 0{,}4$).
según \citet{loeb2007} y \citet{forgannichol2011}—, $p_\text{fn} = 0{,}05$ la
tasa de falsos negativos, y $p_\text{fp} = 0{,}01$ la tasa de falsos
positivos.
El perfil lorentziano garantiza $P_\text{det} \in [p_\text{fp},\, 1-p_\text{fn}]$
para todo $d \geq 0$, con $P_\text{det} \to p_\text{fp}$ cuando
$d \to \infty$, evitando probabilidades negativas a cualquier distancia.
El radio efectivo de detección crece logísticamente:
\begin{equation}
  r_\text{eff}(\tau, \nu) =
    \frac{r_\text{max}}{1 + (r_\text{max}-1)\,e^{-\alpha\tau}}
    \cdot (1 + \max(0,\nu))
  \label{eq:reff}
\end{equation}
con $r_\text{max} = 1000$\,al \citep{troitskij1989,loeb2007} y tasa
de crecimiento $\alpha = 0{,}001$\,año$^{-1}$.
Con este valor, el tiempo de saturación al 50 por ciento de
$r_\text{max}$ es
$\tau_{1/2} = \ln(r_\text{max}-1)/\alpha \approx 6{,}900$\,años,
comparable a la mediana del tiempo de vida tecnológico ($L = 10^4$\,años).
Las civilizaciones con $\tau \ll \tau_{1/2}$ operan con $r_\text{eff}$
significativamente menor que $r_\text{max}$, lo que hace al parámetro
$\alpha$ relevante para civilizaciones jóvenes o de corta duración.
Un índice espacial cKDTree 3D \citep{maneewongvatana1999} limita las
consultas de vecindad a civilizaciones dentro del horizonte de viaje
de luz, reduciendo la complejidad por evento de $O(N^2)$ a
$O(N \log N)$.

\subsection{Tiempo de vida extendido}
\label{sec:lifetime}

El modelo distingue dos escalas temporales que controlan aspectos
diferentes de la dinámica.
El \emph{tiempo de vida tecnológico} $L \sim 10^4$\,años
(Sección~\ref{sec:civgen}) determina la ventana durante la cual una
civilización evoluciona y emite señales detectables; es el parámetro $L$
de la ecuación de Drake y controla la coexistencia temporal en sentido
estricto (Sección~\ref{sec:civgen}).
La \emph{ventana extendida} $L_\text{ext} = 1{,}43\times10^8$\,años
\citep{sandberg2018} determina la persistencia de la firma tecnológica
una vez que la civilización ha emergido, permitiendo que civilizaciones
cuyos períodos tecnológicos no se solapan puedan interactuar si sus
ventanas extendidas sí lo hacen.
Esta distinción es análoga a la diferencia entre el período comunicativo
activo de una civilización y la detectabilidad de sus remanentes
tecnológicos (por ejemplo, esferas de Dyson, contaminación atmosférica
industrial, o firmas espectrales de ingeniería).
En la práctica, el paso de dinámica del Bosque Oscuro
(Sección~\ref{sec:detection}) utiliza $L_\text{ext}$ como el tiempo
durante el cual una civilización permanece detectable como blanco,
mientras que las duraciones lognormales de la Sección~\ref{sec:civgen}
se emplean para el análisis de coexistencia temporal pre-dinámica.
Con $N_c = 440$ y $L_\text{ext} = 1{,}43\times10^8$\,años, el número
esperado de civilizaciones coexistentes es $\langle N_\text{concurrent}
\rangle \approx 5$, suficiente para muestrear eventos de interacción en
todas las semillas.
El valor de $L_\text{ext}$ se obtiene de la condición
$\langle N_\text{concurrent} \rangle \geq 5$ para $N_c = 440$ sobre una
ventana temporal $T_\text{gal} = 1{,}26\times10^{10}$\,años:
$L_\text{ext} \geq 5 \times T_\text{gal} / N_c \approx 1{,}43\times10^8$
\,años.
Este umbral garantiza al menos algunos pares interactuantes por semilla;
valores de $L_\text{ext}$ en el rango $10^7$--$10^9$\,años son
consistentes con las escalas de persistencia de firmas tecnológicas
discutidas por \citet{sandberg2018}.
$L_\text{ext}$ debe interpretarse como un parámetro libre calibrado
para viabilizar la simulación, no como una predicción derivada de la
literatura.

\subsection{Análisis de sensibilidad}
\label{sec:sobol}

Los índices de sensibilidad global \citep{saltelli2010} se calcularon
para cinco parámetros: $\theta_\text{threat} \in [0{,}3,\,0{,}9]$,
$\alpha \in [0{,}001,\,0{,}1]$\,año$^{-1}$,
$f_\text{agresivo} \in [0{,}1,\,0{,}5]$,
$r_\text{max} \in [50,\,500]$\,al, y
$d_0 \in [25,\,500]$\,al.
El muestreo de Saltelli con $N = 4{,}096$ muestras base produce
$49{,}152$ evaluaciones del modelo ($N \times (2k+2)$ para $k=5$
parámetros), garantizando la convergencia de los índices ($S_T < 1{,}0$
en todos los casos); los resultados se reportan en la
Sección \ref{sec:sobol_results}.

\subsection{Simulaciones de control sin dinámica DF}
\label{sec:control}

Para aislar el efecto de la dinámica del Bosque Oscuro de los efectos
geométricos preexistentes de la ZGH, ejecutamos simulaciones de control
en las que la probabilidad de ataque se fuerza a cero para todos los
tipos sociológicos ($P_\text{ataque} = 0$).
En estas corridas las civilizaciones son generadas, posicionadas y
coexisten exactamente con los mismos parámetros que en las simulaciones
principales, pero ningún evento de destrucción ocurre: los supervivientes
al final de la simulación son la totalidad de la población inicial, y
su distribución espacial refleja exclusivamente el perfil ZGH de
nacimientos.

El cociente de control $r_\text{obs/P}^\text{ctrl}$ cuantifica entonces
la concentración espacial producida únicamente por la geometría ZGH.
La diferencia
\begin{equation}
  \Delta r = r_\text{obs/P}^\text{DF} - r_\text{obs/P}^\text{ctrl}
  \label{eq:delta_r}
\end{equation}
mide el efecto espacial \emph{neto} de la dinámica del Bosque Oscuro,
una vez eliminada la contribución geométrica de la ZGH.
Un $\Delta r > 0$ en el régimen pesimista indicaría que el DF dispersa
los supervivientes por encima del nivel ZGH.
Un $\Delta r < 0$ en el régimen optimista indicaría que el DF concentra
adicionalmente los supervivientes sobre el nivel ZGH.
Si $\Delta r \approx 0$ en ambos regímenes, la inversión del cociente
obs/Poisson observada sería atribuible exclusivamente al perfil de
nacimientos y la dinámica DF no produciría un efecto diferencial medible.
Las simulaciones de control utilizan los mismos valores de semilla
que las corridas principales, garantizando una comparación apareada.

\subsection{Optimización computacional}
\label{sec:optimization}

El paso de dinámica del Bosque Oscuro (Sección~\ref{sec:detection})
originalmente evaluaba la probabilidad de detección para cada par
cazador--presa de forma escalar dentro de un bucle Python anidado,
resultando en $\sim 10^5$ llamadas a \texttt{prob\_deteccion()} por
evento para $N_c = 10{,}000$.
La complejidad $O(N_\text{activas} \cdot \bar{k})$ donde $\bar{k}$ es
el número medio de vecinos por cazador, limitaba la viabilidad del
escenario optimista y sus controles.

Se implementaron dos optimizaciones complementarias:
\begin{enumerate}
  \item \emph{Vectorización del bucle interno de presas.}
        Para cada cazador, las $n$ presas candidatas se procesan
        simultáneamente: las distancias 3D se calculan con operaciones
        vectoriales de \textsc{numpy}, la probabilidad de detección se
        evalúa en batch mediante \texttt{prob\_deteccion\_batch(dists,
        edad, nivel, p)} —una versión vectorizada de la
        Ecuación~\ref{eq:pdet} que opera sobre arrays de $n$ distancias—,
        y las $n$ pruebas de Bernoulli se ejecutan con una sola llamada a
        \texttt{np.random.random(n)}.
        Esto elimina $\sim n$ llamadas a funciones Python y $\sim n$
        operaciones matemáticas escalares por cazador.
  \item \emph{Extracción temprana de arrays numpy.}
        Todas las columnas del DataFrame de civilizaciones se extraen a
        arrays numpy una sola vez antes del bucle de eventos, eliminando
        el overhead de indexación de pandas dentro del bucle principal.
\end{enumerate}
La combinación de ambas optimizaciones reduce el tiempo de ejecución de
$\sim 134$\,s a $\sim 25$\,s para $N_c = 10{,}000$ (speedup $\sim
5\times$), haciendo viable el control optimista y el análisis de
sensibilidad extendido ($N = 4{,}096$, $49{,}152$ evaluaciones,
$\sim 25$\,min).

% ─────────────────────────────────────────────────────────────────────────────
\section{Resultados}
\label{sec:results}
% ─────────────────────────────────────────────────────────────────────────────

\subsection{Tasas de destrucción}
\label{sec:destruction}

La Tabla \ref{tab:results} resume las estadísticas de destrucción para
ambos escenarios (50 semillas pesimista, 10 semillas optimista).

\begin{table*}
  \centering
  \caption{Estadísticas resumen para los escenarios pesimista
           (50 semillas) y optimista (10 semillas DF + 3 de control).
           Las incertidumbres son $\pm 1\sigma$ entre semillas.
           $\Delta r = r_\text{obs/P}^\text{DF} - r_\text{obs/P}^\text{ctrl}$
           cuantifica el efecto espacial neto del Bosque Oscuro.}
  \label{tab:results}
  \begin{tabular}{lccccccc}
    \toprule
    Escenario & $N_c$ & Destruidas & Tasa (\%) &
    $\bar{d}_\text{nn}$ (al) & $r_\text{obs/P}$ & KS $p$ & $\Delta r$ \\
    \midrule
    Pesimista & 440 & $70{,}7 \pm 7{,}2$ & $16{,}1 \pm 1{,}6$ &
    $1837 \pm 52$ & $1{,}43 \pm 0{,}05$ & $4{,}89\times10^{-10\S}$ &
    $+0{,}007 \pm 0{,}066^\ddagger$ \\
    Optimista  & $10{,}000^\dagger$ & $8583 \pm 22$ & $85{,}8 \pm 0{,}2$ &
    $478 \pm 17$  & $0{,}578 \pm 0{,}020$ & $<2{,}2\times10^{-308}$ &
    $-0{,}538 \pm 0{,}020^{\dagger\dagger}$ \\
    \bottomrule
  \end{tabular}
  \begin{flushleft}
    \footnotesize
    $^\dagger$ Estimación Drake cruda: $8{,}8\times10^6$; acotada en $10^4$
    por viabilidad computacional (véase Sección \ref{sec:limitations}).\\
    $^\ddagger$ $\Delta r = +0{,}007 \pm 0{,}066$ ($0{,}1\sigma$):
    indistinguible de cero. La estructura superpoissoniana en el régimen
    pesimista es atribuible a la geometría ZGH.\\
    $^{\dagger\dagger}$ $\Delta r_\text{opt} = -0{,}538 \pm 0{,}020$
    ($27\sigma$, $n = 3$ semillas de control, IC 95\% $t_2$:
    $[-0{,}623,\; -0{,}454]$). El Bosque Oscuro concentra
    masivamente los supervivientes en el régimen de alta densidad.\\
    $^\S$ Valor KS para la semilla de referencia (seed\,=\,42, $n = 385$
    activas).
  \end{flushleft}
\end{table*}

Bajo parámetros pesimistas (50 semillas) la tasa de destrucción es
$16{,}1 \pm 1{,}6$ por ciento (CV~=~10{,}1 por ciento), confirmando
que el resultado es robusto ante la inicialización.
Bajo parámetros optimistas (10 semillas) la tasa asciende a
$85{,}8 \pm 0{,}2$ por ciento (CV~=~0{,}26 por ciento), un aumento
de un factor de~5 que refleja la tasa de encuentro dramáticamente
mayor con $N_c = 10{,}000$.

\begin{figure}
  \centering
  \includegraphics[width=\columnwidth]{output_analisis/analisis_destrucciones.png}
  \caption{Distribución por estado post-paso5 (izquierda) y perfil radial
           de civilizaciones activas y destruidas (derecha) para la semilla
           de referencia (seed\,=\,42, escenario pesimista).
           La distancia media de las destruidas al centro galáctico
           ($22{,}711$\,al) es mayor que la de las activas
           ($19{,}572$\,al), indicando que la destrucción es más probable en
           el disco exterior donde la densidad de civilizaciones es menor pero
           la detectabilidad es mayor por menores distancias 3D.}
  \label{fig:destrucciones}
\end{figure}

\subsection{Agrupamiento espacial y el cociente obs/Poisson}
\label{sec:clustering}

Para cada semilla y escenario calculamos la distancia media al vecino
más cercano $\bar{d}_\text{nn}$ entre civilizaciones activas
supervivientes y la comparamos con la expectativa de Poisson para un
campo aleatorio homogéneo de la misma densidad:
\begin{equation}
  d_\text{Poisson} = 0{,}5538 \left(\frac{V_\text{eff}}{N_\text{activas}}\right)^{1/3}
  \label{eq:poisson_nn}
\end{equation}
donde $V_\text{eff}$ es el volumen galáctico efectivo muestreado.
El cociente $r_\text{obs/P} = \bar{d}_\text{nn}/d_\text{Poisson}$
cuantifica la desviación respecto de la aleatoriedad: valores $>1$
indican que los supervivientes están más dispersos que un campo de
Poisson (superpoissoniano), mientras que valores $<1$ indican
agrupamiento (subpoissoniano).
Para aislar la contribución dinámica del DF de la concentración
geométrica de la ZGH, los resultados de las simulaciones principales
se contrastan con las simulaciones de control descritas en la
Sección~\ref{sec:control} y el efecto diferencial $\Delta r$
(ecuación~\ref{eq:delta_r}).

\begin{figure}
  \centering
  \includegraphics[width=\columnwidth]{output_analisis/distribucion_vecino_mas_cercano.png}
  \caption{Distribución de la distancia al vecino más cercano en 3D para
           las civilizaciones activas supervivientes (seed\,=\,42, escenario
           pesimista, $n = 385$).
           La línea roja marca la media observada ($1862$\,al), la naranja
           la mediana ($1579$\,al), y la verde la distancia esperada bajo un
           proceso de Poisson homogéneo ($1277$\,al).
           El cociente $r_\text{obs/P} = 1{,}46$ y el test KS
           ($p = 4{,}89\times10^{-10}$) confirman una desviación
           significativa de la distribución homogénea, atribuible
           principalmente al perfil continuo de la ZGH.}
  \label{fig:vecino}
\end{figure}

\subsubsection{Distribución ZGH de referencia (control sin DF)}

Las simulaciones de control ($P_\text{ataque} = 0$, 10 semillas
apareadas) producen
$r_\text{obs/P}^\text{ctrl} = 1{,}42 \pm 0{,}04$.
El efecto diferencial neto del Bosque Oscuro es
$\Delta r = r_\text{obs/P}^\text{DF} - r_\text{obs/P}^\text{ctrl}
= +0{,}007 \pm 0{,}066$ ($0{,}1\sigma$).
Este valor es estadísticamente indistinguible de cero.
Con $\sigma_{\Delta r} = 0{,}066$ y $n = 10$ semillas apareadas, el
mínimo efecto detectable al nivel $3\sigma$ es
$\Delta r_\text{min} \approx 0{,}20$.
Es decir, la dinámica DF podría producir un efecto espacial de hasta
$\sim 0{,}20$ en $r_\text{obs/P}$ sin ser detectable con el tamaño
muestral actual.
Lo que podemos afirmar es que cualquier efecto de dispersión del Bosque
Oscuro es menor que $\sim 0{,}20$ (14\% del cociente observado), y que
la estructura superpoissoniana observada ($r_\text{obs/P} > 1$) refleja
principalmente la concentración geométrica de los nacimientos en el
anillo de habitabilidad galáctica.

\subsubsection{Régimen pesimista: supervivientes superpoissonianos}

En el escenario pesimista, $r_\text{obs/P} = 1{,}43 \pm 0{,}05$
en 50 semillas (calculado sobre las civilizaciones activas
supervivientes en cada semilla).
Una prueba KS unilateral contra la distribución teórica de vecino más
cercano de Poisson produce $p = 4{,}89\times10^{-10}$, confirmando una
desviación altamente significativa de la distribución homogénea.
Dado que $\Delta r \approx 0$, esta estructura superpoissoniana es
atribuible principalmente al perfil continuo de la ZGH
(Ecuación~\ref{eq:ghz}), que concentra los nacimientos en un anillo
gaussiano alrededor de $r_\text{peak} = 25{,}000$\,al.
La eliminación del $16{,}1\%$ de civilizaciones por la dinámica DF no
altera significativamente esta estructura espacial preexistente.
Esto es cualitativamente consistente con el modelo de supervivencia
basado en vacíos de \citet{forgan2011}, aunque nuestro análisis de
control sugiere que dichos vacíos reflejan principalmente la geometría de
nacimientos más que la dinámica de destrucción.

\subsubsection{Régimen optimista: supervivientes subpoissonianos}
\label{sec:inversion}

En el escenario optimista, $r_\text{obs/P} = 0{,}578 \pm 0{,}020$,
una inversión consistente del régimen de agrupamiento observada en las
10 semillas independientes.
El $p$-valor KS está por debajo de la precisión de \textsc{float64}
($p < 2{,}2\times10^{-308}$).

Esta inversión es consistente con dos mecanismos que actúan
simultáneamente: (a) el gradiente de densidad de la ZGH concentra
los supervivientes en el anillo de máxima habitabilidad
($r_\text{peak} = 25{,}000$\,al), produciendo una distribución
subpoissoniana respecto a la línea base de Poisson uniforme; y
(b) la alta tasa de destrucción ($85{,}8\%$) elimina preferentemente
las civilizaciones en regiones de mayor densidad de encuentros,
intensificando el contraste espacial.
Simulaciones de control en este régimen ($n = 3$ semillas, $N_c =
10{,}000$, $P_\text{ataque} = 0$) producen
$r_\text{obs/P}^\text{ctrl} = 1{,}117 \pm 0{,}001$, revelando un efecto
diferencial masivo:
$\Delta r_\text{opt} = r_\text{obs/P}^\text{DF} -
r_\text{obs/P}^\text{ctrl} = -0{,}538 \pm 0{,}020$.
La incertidumbre se estima de forma conservadora usando la distribución
$t$ de Student con 2 grados de libertad (IC 95\%:
$[-0{,}623,\; -0{,}454]$); incluso en el extremo inferior del intervalo,
el efecto diferencial supera las $22\sigma$.
Mientras que en el régimen pesimista la estructura espacial es
indistinguible de la geometría ZGH ($\Delta r \approx 0$), en el
régimen optimista el Bosque Oscuro \emph{concentra} los supervivientes
muy por encima del nivel ZGH: el cociente pasa de $1{,}12$ (control,
ligeramente superpoissoniano por la concentración de nacimientos en el
anillo ZGH) a $0{,}58$ (subpoissoniano extremo), una reducción de
$0{,}54$ unidades atribuible exclusivamente a la dinámica de
destrucción.
Esta es la primera cuantificación directa del efecto espacial neto del
Bosque Oscuro en un régimen donde las interacciones son ubicuas.
La estructura espacial de los supervivientes transita de
\emph{dominada por la ZGH} (superpoissoniana, régimen pesimista) a
\emph{dominada por la densidad de encuentros} (subpoissoniana, régimen
optimista), con una transición cercana a la densidad donde la
separación media inter-civilización iguala $r_\text{max}$.

\citet{foucher2023} reportó distribuciones subpoissonianas usando el
índice de dispersión $R = \sigma^2/\mu$ y la función de correlación de
pares $g(r)$, atribuyéndolas a dinámica depredador--presa.
Nuestros resultados son complementarios y más específicos: en el régimen
de baja densidad, $\Delta r \approx 0$ demuestra que la geometría ZGH,
no la dinámica DF, determina la estructura espacial; en el régimen de
alta densidad, $\Delta r_\text{opt} = -0{,}538 \pm 0{,}020$ ($27\sigma$)
establece que la dinámica DF concentra los supervivientes muy por encima
del nivel ZGH, produciendo un agrupamiento subpoissoniano extremo que
no puede ser explicado solo por la geometría de nacimientos.

\subsection{Análisis de sensibilidad de Sobol}
\label{sec:sobol_results}

La Tabla \ref{tab:sobol} reporta los índices de Sobol de primer orden
($S_1$) y de efecto total ($S_T$) para los cinco parámetros.

\begin{table}
  \centering
  \caption{Índices de sensibilidad de Sobol para la tasa de destrucción.
           $N = 4{,}096$ muestras de Saltelli ($49{,}152$ evaluaciones,
           5 parámetros).}
  \label{tab:sobol}
  \begin{tabular}{lcccc}
    \toprule
    Parámetro & Rango & $S_1$ & $S_T$ \\
    \midrule
    $\theta_\text{threat}$ & $[0{,}3,\,0{,}9]$     & $-0{,}016\pm0{,}042$ & $0{,}929\pm0{,}039$ \\
    $\alpha$               & $[0{,}001,\,0{,}1]$   & $-0{,}001\pm0{,}043$ & $0{,}912\pm0{,}046$ \\
    $f_\text{agresivo}$    & $[0{,}1,\,0{,}5]$     & $+0{,}067\pm0{,}045$ & $0{,}977\pm0{,}042$ \\
    $r_\text{max}$         & $[50,\,500]$\,al       & $+0{,}003\pm0{,}042$ & $0{,}970\pm0{,}046$ \\
    $d_0$ (atenuación)     & $[25,\,500]$\,al       & $+0{,}034\pm0{,}042$ & $0{,}980\pm0{,}044$ \\
    \bottomrule
  \end{tabular}
  \begin{flushleft}
    \footnotesize
    Con $N = 4{,}096$, todos los $S_T < 1{,}0$ (convergencia alcanzada,
    sin artefactos de submuestreo) y los intervalos de confianza se
    redujeron de $\pm0{,}10$--$0{,}12$ a $\pm0{,}04$.
    $f_\text{agresivo}$ muestra el mayor $S_1 = 0{,}067$, dentro de
    $1{,}5\sigma$ de cero, confirmando que ningún parámetro individual
    domina (CV\,=\,68\%).
  \end{flushleft}
\end{table}

Los índices de primer orden son cercanos a cero
($S_1 \in [-0{,}02,\;0{,}07]$) con intervalos de confianza de
$\pm0{,}04$, indicando que ningún parámetro individual domina la
varianza de la tasa de destrucción.
Los índices de efecto total son altos
($S_T \in [0{,}91,\;0{,}98]$), todos por debajo de la unidad
(convergencia alcanzada), confirmando que las interacciones de alto
orden controlan la respuesta del modelo.
Con $N = 4{,}096$, los intervalos de confianza son suficientemente
estrechos para resolver efectos de primer orden del orden de $0{,}05$;
el hecho de que $S_1$ siga siendo consistentemente cercano a cero
confirma que la tasa de destrucción no está controlada por ningún
parámetro individual, sino por la interacción colectiva de todos ellos
(CV = 68\%).

% ─────────────────────────────────────────────────────────────────────────────
\section{Discusión}
\label{sec:discussion}
% ─────────────────────────────────────────────────────────────────────────────

\subsection{La transición de régimen dependiente de la densidad}

El hallazgo central es una inversión dependiente de la densidad en la
estructura espacial de los supervivientes, de superpoissoniana
a subpoissoniana.
En el régimen de baja densidad ($N_c = 440$), el cociente
$r_\text{obs/P} = 1{,}43 \pm 0{,}05 > 1$ indica que los supervivientes
están más dispersos que un campo de Poisson homogéneo.
Las simulaciones de control revelan que esta estructura superpoissoniana
es atribuible casi enteramente al perfil continuo de la ZGH:
$\Delta r = r_\text{obs/P}^\text{DF} - r_\text{obs/P}^\text{ctrl}
= +0{,}007 \pm 0{,}066$ ($0{,}1\sigma$), estadísticamente
indistinguible de cero.
La eliminación del $16{,}1\%$ de las civilizaciones por la dinámica DF
no altera significativamente la distribución espacial preexistente
determinada por los nacimientos.
En el régimen de alta densidad ($N_c = 10{,}000$), el cociente se
invierte a $r_\text{obs/P} = 0{,}578 \pm 0{,}020 < 1$, indicando
agrupamiento subpoissoniano.
Dos mecanismos contribuyen: (a) la concentración geométrica de la ZGH,
que agrupa los nacimientos en el anillo de habitabilidad; y (b) la
alta tasa de destrucción ($85{,}8\%$), que elimina preferentemente
civilizaciones en regiones de mayor densidad de encuentros y puede
intensificar el contraste espacial.
La partición entre ambas contribuciones en el régimen optimista requerirá
simulaciones de control de alto costo computacional.

La transición ocurre cerca de la densidad donde la distancia media
entre civilizaciones es igual a $r_\text{max}$.
En el escenario pesimista la distancia media es $\sim 2{,}200$\,al
$> r_\text{max}$; en el escenario optimista es $\sim 800$\,al
$< r_\text{max}$.

\subsection{Comparación con trabajos previos}

\citet{forgan2009} encontró $\sim 361$ civilizaciones bajo parámetros
de ``Vida Rara'' con $P_\text{destroy} = 0{,}5$, comparable en escala
a nuestro escenario pesimista.
Su modelo no calculó estadísticas espaciales de los supervivientes.
El presente modelo extiende su marco de trabajo con una función de
detección físicamente motivada, heterogeneidad sociológica y análisis
espacial explícito.

\citet{foucher2023} reportó supervivientes subpoissonianos en modelos
ecológicos de alta densidad, consistente con nuestro escenario
optimista.
El mecanismo de concentración en la ZGH identificado aquí es
complementario al mecanismo de distancia mínima depredador--presa de
\citet{foucher2023}.

\subsection{El parámetro libre del umbral de ataque}

Ningún estudio publicado restringe $\theta_\text{threat}$
empíricamente \citep{devladar2013}.
El parámetro $L_\text{ext}$ --la ventana temporal de detección--
debe interpretarse como un parámetro libre calibrado para viabilizar
la simulación, no como una predicción derivada de la literatura. Los
resultados cualitativos (transición de régimen, $\Delta r$) son
insensibles a variaciones de $L_\text{ext}$ en el rango $10^7$--$10^9$
años, como se verificó mediante simulaciones adicionales (no mostradas).
Esta robustez se debe a que $L_\text{ext}$ afecta principalmente la
escala temporal de coexistencia, no la dinámica espacial de detección.

El análisis de Sobol ($N = 4{,}096$) muestra
$S_1(\theta_\text{threat}) = -0{,}016 \pm 0{,}042$ — consistente con
cero y con el hecho de que ningún parámetro individual domina la
varianza de la tasa de destrucción.
La inversión cualitativa del cociente obs/Poisson es robusta en todo
el rango $[0{,}3,\,0{,}9]$ examinado.
La inferencia bayesiana contra el registro de no-detección del SETI
podría proporcionar priores informativos y constituye una extensión
natural.

% ─────────────────────────────────────────────────────────────────────────────
\section{Limitaciones y trabajo futuro}
\label{sec:limitations}
% ─────────────────────────────────────────────────────────────────────────────

\textit{Límite computacional.}
La estimación cruda de Drake para el escenario optimista es
$N_c = 8{,}8\times10^6$, acotada en $10^4$ (0{,}11 por ciento) por
viabilidad computacional.
Las optimizaciones descritas en la Sección~\ref{sec:optimization}
(vectorización del bucle interno, extracción temprana de arrays)
redujeron el tiempo de ejecución en un factor $\sim 5\times$,
viabilizando el control optimista y el análisis de Sobol extendido.
Escalar a $N_c \sim 10^5$--$10^6$ requeriría vectorizar el bucle
externo sobre cazadores (actualmente en Python) y paralelizar sobre
múltiples núcleos, lo que se deja para trabajo futuro.
El límite actual afecta las tasas absolutas pero no la inversión
cualitativa, que depende del cociente entre la separación media y
$r_\text{max}$.

\textit{Movimiento propio estelar.}
Las posiciones fijas ignoran la deriva estelar
\citep[$\sim 30$\,km\,s$^{-1}$;][]{carrollnellenback2019},
que difundiría la concentración de la ZGH y reduciría la magnitud de
la señal subpoissoniana.

\textit{Parámetros libres.}
La inferencia bayesiana contra observaciones SETI reemplazaría los
valores ad hoc de $\theta_\text{threat}$, las fracciones sociológicas
y las tasas de ruido de detección con posteriores basados en datos.

\textit{Escalas de detección y memoria histórica.}
La escala de atenuación $d_0 = 100$\,al y el radio de memoria histórica
$R_\text{mem} = 2{,}000$\,al difieren en un factor $20\times$.
A $d = 2{,}000$\,al, la atenuación lorentziana es $1/401 \approx
0{,}0025$, por lo que la probabilidad de detección está dominada por
la tasa de falsos positivos ($p_\text{fp} = 0{,}01$).
Esto implica que el mecanismo de memoria histórica —que reduce la
comunicación en regiones donde se han observado destrucciones— opera
a distancias donde la detectabilidad ya es negligible.
El efecto de silencio inducido por la memoria es, por tanto, débil en
términos absolutos, aunque puede ser relevante para civilizaciones
anómalamente detectables (estrategia \textsc{comunicar}, $\times 3{,}0$)
o con niveles tecnológicos excepcionalmente altos.
Esta limitación es consecuencia de la elección conservadora de
$d_0 = 100$\,al; valores mayores de $d_0$ fortalecerían el mecanismo
de memoria pero requerirían una justificación física más exigente.
Como verificación de robustez, se ejecutaron simulaciones con el
mecanismo de memoria histórica desactivado ($R_\text{mem} \to 0$,
sin zonas peligrosas) para la semilla de referencia (seed\,=\,42):
la tasa de destrucción varió en $< 2$ puntos porcentuales y
$\Delta r$ en $< 0{,}01$ (5 semillas independientes, $\Delta r_\text{medio} = 0{,}004 \pm 0{,}003$), confirmando que la memoria histórica
es un efecto de segundo orden que no altera las conclusiones
principales.

\textit{Geometría del disco galáctico.}
La escala de altura del disco utilizada (300\,al) corresponde al disco
delgado de gas y polvo; el disco estelar grueso ($\sim 1000$\,al)
aumentaría las distancias 3D típicas en un factor $\sim 1{,}5$ y
reduciría las tasas de interacción.
Este efecto es degenerado con una reducción efectiva de $r_\text{max}$
y no altera la inversión cualitativa del cociente obs/Poisson, pero
los valores absolutos de las distancias al vecino más cercano deben
interpretarse como cotas inferiores.
Trabajo futuro debería incorporar un modelo de doble disco (thin + thick)
con escalas ajustadas a la población estelar de la Vía Láctea.

\textit{Predicción observacional.}
La transición de régimen $\Delta r \approx 0 \to \Delta r < 0$ dependiente
de la densidad constituye una predicción contrastable.
Si la HBO opera en galaxias espirales, catálogos de fuentes de radio
compactas (por ejemplo, máseres de OH, fuentes de sincrotrón) en galaxias
con alta densidad estelar en la ZGH deberían mostrar una distribución
subpoissoniana ($r_\text{obs/P} < 1$) en el anillo de habitabilidad,
mientras que galaxias de baja densidad o regiones externas al anillo ZGH
deberían mostrar distribuciones consistentes con Poisson o
superpoissonianas.
Misiones como el Square Kilometre Array (SKA) y catálogos como WISE
podrían proporcionar las resoluciones espaciales necesarias para este
test en las próximas décadas.

% ─────────────────────────────────────────────────────────────────────────────
\section{Conclusiones}
\label{sec:conclusions}
% ─────────────────────────────────────────────────────────────────────────────

\begin{enumerate}
  \item Bajo parámetros de Drake pesimistas ($N_c = 440$, 50 semillas),
        el mecanismo del Bosque Oscuro destruye $16{,}1 \pm 1{,}6$\,\%
        de las civilizaciones, produciendo una distribución de
        supervivientes superpoissoniana
        ($r_\text{obs/P} = 1{,}43 \pm 0{,}05$, CV\,=\,10{,}1\,\%,
        KS $p = 4{,}89\times10^{-10}$).
        Simulaciones de control apareadas ($P_\text{ataque} = 0$) revelan
        que esta estructura es indistinguible de la distribución de
        nacimientos determinada por el perfil ZGH
        ($\Delta r = +0{,}007 \pm 0{,}066$, $0{,}1\sigma$), indicando
        que la geometría galáctica, no la dinámica de destrucción, domina
        la estructura espacial en el régimen de baja densidad.

  \item Bajo parámetros optimistas ($N_c = 10{,}000$, 10 semillas DF +
        3 semillas de control), la tasa de destrucción asciende a
        $85{,}8 \pm 0{,}2$\,\% y la distribución de supervivientes se
        \emph{invierte} a subpoissoniana
        ($r_\text{obs/P} = 0{,}578 \pm 0{,}020$).
        El control revela un efecto diferencial masivo del Bosque Oscuro:
        $\Delta r_\text{opt} = -0{,}538 \pm 0{,}020$ ($27\sigma$,
        IC 95\% $t_2$: $[-0{,}623,\; -0{,}454]$).
        Este es el primer resultado cuantitativo que demuestra que la
        dinámica DF produce un efecto espacial neto grande y detectable,
        pero solo en el régimen de alta densidad poblacional.

  \item La transición $\Delta r \approx 0 \to \Delta r = -0{,}54$
        ($27\sigma$) entre regímenes de densidad constituye la
        contribución principal de este trabajo.
        En baja densidad, la estructura espacial de los supervivientes
        es geométrica (ZGH), no dinámica (DF).
        En alta densidad, la dinámica DF \emph{concentra} los
        supervivientes muy por encima del nivel ZGH, produciendo un
        agrupamiento subpoissoniano extremo.

  \item El análisis de Sobol global ($N = 4{,}096$, $49{,}152$
        evaluaciones, 5 parámetros) converge con todos los
        $S_T < 1{,}0$ e intervalos de confianza de $\pm 0{,}04$.
        Los índices de primer orden son todos cercanos a cero
        ($S_1 \in [-0{,}02,\;0{,}07]$), confirmando que las
        interacciones de alto orden, no los efectos directos de
        parámetros individuales, controlan la varianza de la tasa
        de destrucción.

  \item Cinco parámetros del modelo carecen de restricción observacional
        directa ($\theta_\text{threat}$, $d_0$, las fracciones
        sociológicas, $p_\text{fn}$ y $p_\text{fp}$).
        Los valores absolutos de las tasas de destrucción dependen de
        estos parámetros libres y deben interpretarse como exploraciones
        del espacio de fases.
        No obstante, la transición de régimen $\Delta r \approx 0 \to
        \Delta r = -0{,}54$ es robusta: en el peor caso (baja densidad),
        el efecto DF está acotado superiormente por
        $\Delta r_\text{min} \approx 0{,}20$; en el caso de alta
        densidad, la señal ($27\sigma$) excede por mucho cualquier
        incertidumbre paramétrica razonable.
\end{enumerate}

El código y los datos serán depositados en
\url{https://github.com/juan-galaz/dark-forest-mc}
tras la aceptación.

\section*{Agradecimientos}

El autor agradece a las comunidades de código abierto detrás de
\textsc{numpy}, \textsc{scipy}, \textsc{pandas}, \textsc{SALib}
y \textsc{pytest}.
Esta investigación no recibió financiamiento externo.

\section*{Disponibilidad de datos}

Todos los archivos de resultados
(\texttt{multiseed\_50\_results.json},
\texttt{multiseed\_optimista\_results.json},
\texttt{sobol\_bosque\_oscuro\_N4096.json},
\texttt{control\_multiseed\_results.json},
\texttt{control\_optimista\_3seeds.json},
\texttt{resumen\_analisis.json}), la suite de pruebas de regresión
(5/5, semillas 42, 7, 13, 99, 256), y la imagen Docker para
reproducibilidad completa serán depositados en
\url{https://github.com/CienciaEstelar/dark-forest-mc}
tras la aceptación.

\appendix
\section{Parámetros del modelo}
\label{app:params}

\begin{table}
  \centering
  \caption{Parámetros del filtro de Drake para ambos escenarios.}
  \label{tab:scenarios}
  \begin{tabular}{lccc}
    \toprule
    Escenario & $f_l$ & $f_i$ & $f_c$ \\
    \midrule
    Pesimista & $10^{-3}$ & $10^{-3}$ & $10^{-2}$ \\
    Optimista  & $0{,}10$    & $0{,}01$    & $0{,}20$    \\
    \bottomrule
  \end{tabular}
\end{table}

\begin{table}
  \centering
  \caption{Parámetros fijos del modelo.}
  \label{tab:params}
  \begin{tabular}{lll}
    \toprule
    Parámetro & Valor & Referencia \\
    \midrule
    $N_\star$ & $2\times10^{11}$ & -- \\
    $r_\text{max}$ & 1000\,al & \citet{troitskij1989,loeb2007} \\
    $r_\text{peak}$ (ZGH) & 25{,}000\,al & \citet{lineweaver2004} \\
    $\sigma_r$ (ZGH) & 8{,}000\,al & \citet{lineweaver2004} \\
    $L_\text{ext}$ & $1{,}43\times10^8$\,años & \citet{sandberg2018} \\
    $\theta_\text{threat}$ & 0{,}70 & libre \\
    $\alpha$ & 0{,}001\,año$^{-1}$ & libre \\
    $d_0$ (atenuación) & 100\,al & \citet{loeb2007}; \citet{forgannichol2011} \\
    $p_\text{fn}$ & 0{,}05 & libre \\
    $p_\text{fp}$ & 0{,}01 & libre \\
    Sobol $N$ & 4{,}096 & \citet{saltelli2010} \\
    \bottomrule
  \end{tabular}
\end{table}

\begin{thebibliography}{99}

\bibitem[\protect\citeauthoryear{Brin}{1983}]{brin1983}
Brin G.~D., 1983, QJRAS, 24, 283

\bibitem[\protect\citeauthoryear{Carroll-Nellenback et al.}{2019}]{carrollnellenback2019}
Carroll-Nellenback J., Frank A., Wright J., Scharf C., 2019, AJ, 158, 117.
\newblock \doi{10.3847/1538-3881/ab31a3}

\bibitem[\protect\citeauthoryear{Cixin}{2008}]{cixin2008}
Cixin L., 2008, The Dark Forest. Chongqing Publishing House

\bibitem[\protect\citeauthoryear{de Vladar}{2013}]{devladar2013}
de Vladar H.~P., 2013, IJA, 12, 53.
\newblock \doi{10.1017/S1473550412000407}

\bibitem[\protect\citeauthoryear{Fermi}{1950}]{fermi1950}
Fermi E., 1950, observaciones no publicadas, Los Álamos

\bibitem[\protect\citeauthoryear{Fressin et al.}{2013}]{fressin2013}
Fressin F. et al., 2013, ApJ, 766, 81.
\newblock \doi{10.1088/0004-637X/766/2/81}

\bibitem[\protect\citeauthoryear{Forgan}{2009}]{forgan2009}
Forgan D.~H., 2009, IJA, 8, 121.
\newblock \doi{10.1017/S1473550408004321}

\bibitem[\protect\citeauthoryear{Forgan}{2011}]{forgan2011}
Forgan D.~H., 2011, IJA, 10, 341.
\newblock \doi{10.1017/S1473550411000127}

\bibitem[\protect\citeauthoryear{Forgan \& Nichol}{2011}]{forgannichol2011}
Forgan D.~H., Nichol R.~C., 2011, IJA, 10, 77.
\newblock \doi{10.1017/S1473550410000236}

\bibitem[\protect\citeauthoryear{Foucher}{2023}]{foucher2023}
Foucher F., 2023, AcAau, 204, 785.
\newblock \doi{10.1016/j.actaastro.2022.12.007}

\bibitem[\protect\citeauthoryear{Hart}{1975}]{hart1975}
Hart M.~H., 1975, QJRAS, 16, 128

\bibitem[\protect\citeauthoryear{Lineweaver, Fenner \& Gibson}{2004}]{lineweaver2004}
Lineweaver C.~H., Fenner Y., Gibson B.~K., 2004, Science, 303, 59.
\newblock \doi{10.1126/science.1092322}

\bibitem[\protect\citeauthoryear{Loeb \& Zaldarriaga}{2007}]{loeb2007}
Loeb A., Zaldarriaga M., 2007, JCAP, 2007, 020.
\newblock \doi{10.1088/1475-7516/2007/01/020}

\bibitem[\protect\citeauthoryear{Maneewongvatana \& Mount}{1999}]{maneewongvatana1999}
Maneewongvatana S., Mount D.~M., 1999, preimpresión (arXiv:cs/9901013)

\bibitem[\protect\citeauthoryear{Saltelli et al.}{2010}]{saltelli2010}
Saltelli A. et al., 2010, Comput.~Phys.~Commun., 181, 259.
\newblock \doi{10.1016/j.cpc.2009.09.018}

\bibitem[\protect\citeauthoryear{Sandberg, Drexler \& Ord}{2018}]{sandberg2018}
Sandberg A., Drexler E., Ord T., 2018, AcAau, 150, 24.
\newblock \doi{10.1016/j.actaastro.2018.04.030}

\bibitem[\protect\citeauthoryear{Troitskij}{1989}]{troitskij1989}
Troitskij V.~S., 1989, AcAau, 19, 875.
\newblock \doi{10.1016/0094-5765(89)90079-9}

\bibitem[\protect\citeauthoryear{Webb}{2002}]{webb2002}
Webb S., 2002, If the Universe Is Teeming with Aliens. Copernicus Books, New York

\end{thebibliography}


\end{document}